\documentclass[letterpaper,11pt]{article}
\usepackage[T1]{fontenc}
\usepackage[utf8]{inputenc}
\usepackage{lmodern}
\usepackage[french]{babel}
\usepackage[margin=1in,headheight=28pt]{geometry}
\renewcommand{\familydefault}{\sfdefault}
\usepackage{amsmath,amssymb,amsthm,bm,enumitem}
\usepackage[dvipsnames]{xcolor}
\usepackage[small,compact]{titlesec}
\usepackage{fancyhdr,fancyvrb}
\usepackage[colorlinks=true,allcolors=blue]{hyperref}
\pagestyle{fancy}
\fancyhead[L]{STT-2100 \\ Régression}
\fancyhead[R]{A.-A. Vallée \\ Université Laval}
\fancyfoot[C]{\thepage}
\setlength{\parindent}{0pt}
\setlength{\parskip}{5pt}
\fvset{fontsize=\small,frame=single,framesep=2mm}
\newcommand{\esp}{\mathbb{E}}
\newcommand{\var}{\operatorname{Var}}
\newcommand{\bo}[1]{\bm{#1}}
\newcommand{\N}{\mathcal{N}}
% Définir \avecsolutions avant l'input pour une version corrigée.
\newcommand{\prof}[1]{\ifdefined\avecsolutions\par{\color{blue}\textbf{Solution :} #1}\par\fi}
\begin{document}
\begin{center}{\Large Techniques de régression avancées avec R}\end{center}
Cette fiche accompagne les notes du chapitre 11. Les noms des objets peuvent être adaptés à vos données. Les exemples utilisent des données incluses dans R ou des données simulées; aucun fichier externe n’est requis.

\par\medskip\noindent\begin{minipage}{\linewidth}
\subsection*{Régression polynomiale}
I() donne au carré son sens arithmétique dans une formule. Le modèle reste linéaire en ses coefficients. poly(x, 2) utilise par défaut une base orthogonale : ses coefficients ne sont pas ceux de 1, x, x carré.
\begin{Verbatim}
mod_lin <- lm(dist ~ speed, data = cars)
mod_quad <- lm(dist ~ speed + I(speed^2), data = cars)
anova(mod_lin, mod_quad, test = "F")
summary(mod_quad)
# Conserver les termes de degre inferieur.
plot(cars)
grille <- data.frame(speed = seq(min(cars$speed),
                                max(cars$speed), length.out = 100))
lines(grille$speed, predict(mod_quad, newdata = grille),
      col = "red")
\end{Verbatim}

\end{minipage}\par
\par\medskip\noindent\begin{minipage}{\linewidth}
\subsection*{Régression non linéaire}
Données simulées, avec une graine fixée pour reproduire l'exemple. nls() minimise la somme des carrés sur l'échelle originale. Les valeurs initiales et la convergence doivent être examinées; les tests et IC locaux sont approximatifs.
\begin{Verbatim}
set.seed(2100)
x <- seq(0.1, 3, length.out = 40)
donnees_nl <- data.frame(x = x,
  y = 28 - 34 * exp(-x / 0.6) + rnorm(length(x), sd = 0.5))
mod_nl <- nls(y ~ a + b * exp(-x / theta), data = donnees_nl,
              start = list(a = 28, b = -30, theta = 0.7))
summary(mod_nl)
mod_nl$convInfo
predict(mod_nl, newdata = data.frame(x = c(0.5, 1, 2)))
\end{Verbatim}

\end{minipage}\par
\par\medskip\noindent\begin{minipage}{\linewidth}
\subsection*{Moindres carrés pondérés}
Ici la variance simulée est proportionnelle à x carré et les poids sont donc 1/x carré. Les poids doivent être positifs. Les résidus pondérés sont adaptés au diagnostic de dispersion; des poids estimés ajoutent de l'incertitude.
\begin{Verbatim}
set.seed(2101)
x <- seq(1, 10, length.out = 60)
donnees_w <- data.frame(x = x,
  y = 5 + 2 * x + rnorm(length(x), sd = 0.5 * x))
mod_w <- lm(y ~ x, data = donnees_w, weights = 1 / x^2)
summary(mod_w)
confint(mod_w)
plot(fitted(mod_w), weighted.residuals(mod_w),
     xlab = "Valeurs ajustees", ylab = "Residus ponderes")
abline(h = 0, lty = 2)
\end{Verbatim}

\end{minipage}\par
\par\medskip\noindent\begin{minipage}{\linewidth}
\subsection*{Modéliser une corrélation AR(1)}
Le paquet recommandé nlme permet un ajustement par moindres carrés généralisés. Les données simulées suivent une grille temporelle régulière. L'ordre du temps et la structure de covariance font partie du modèle.
\begin{Verbatim}
if (requireNamespace("nlme", quietly = TRUE)) {
  set.seed(2102)
  temps <- seq_len(100)
  erreur <- as.numeric(arima.sim(list(ar = 0.6), n = 100))
  serie <- data.frame(temps = temps,
                       y = 2 + 0.05 * temps + erreur)
  mod_ar <- nlme::gls(y ~ temps, data = serie,
    correlation = nlme::corAR1(form = ~ temps), method = "REML")
  summary(mod_ar)
  acf(residuals(mod_ar, type = "normalized"))
} else {
  message("L'exemple AR(1) necessite le paquet nlme.")
}
\end{Verbatim}

\end{minipage}\par
\par\medskip\noindent\begin{minipage}{\linewidth}
\subsection*{Aide}
Dans R, utilisez \texttt{?lm} et l’aide des fonctions présentées. 
\end{minipage}\par
\end{document}
